\documentclass[11pt]{article}
\usepackage{amsmath,amssymb,amsthm}
\usepackage{algorithm}
\usepackage{algpseudocode}
\usepackage{fullpage}
\newtheorem{theorem}{Theorem}
\newtheorem{lemma}{Lemma}
\newtheorem{definition}{Definition}
\newcommand{\prox}{\operatorname{prox}}
\newcommand{\argmin}{\operatorname*{arg\,min}}

\begin{document}

\section*{Proximal Gradient Method (PGM)}

\section*{Mathematical Formulation}
Let $f:\mathbb{R}^{d}\rightarrow\mathbb{R}$ be convex and differentiable with $L$‑Lipschitz continuous gradient, i.e.
\[
\|\nabla f(x)-\nabla f(y)\|\le L\|x-y\|\quad\forall x,y\in\mathbb{R}^{d}.
\]
Let $g:\mathbb{R}^{d}\rightarrow\mathbb{R}\cup\{+\infty\}$ be a proper, closed, convex (possibly non‑smooth) function.  
Define the composite objective
\[
F(x)\;:=\;f(x)+g(x).
\]

For a stepsize $\eta>0$, the proximal operator of $g$ is
\[
\prox_{\eta g}(v)\;:=\;\argmin_{u\in\mathbb{R}^{d}}\Big\{ g(u)+\frac{1}{2\eta}\|u-v\|^{2}\Big\}.
\]

\textbf{Update rule:}
\[
\boxed{\,x_{k+1}\;=\;\prox_{\eta g}\!\big(x_{k}-\eta\nabla f(x_{k})\big)\,},\qquad k=0,1,2,\dots
\]

\section*{Pseudocode}
\begin{algorithm}[H]
\caption{Proximal Gradient Method}
\begin{algorithmic}[1]
\Require Initial point $x_{0}\in\mathbb{R}^{d}$, stepsize $\eta\in(0,1/L]$.
\For{$k=0,1,2,\dots$}
    \State Compute gradient $g_{k}\gets \nabla f(x_{k})$.
    \State Take a gradient step $y_{k}\gets x_{k}-\eta g_{k}$.
    \State Apply proximal map $x_{k+1}\gets \prox_{\eta g}(y_{k})$.
\EndFor
\Ensure Approximate solution $x_{K}$ after $K$ iterations.
\end{algorithmic}
\end{algorithm}

\section*{Dry Run Example}
Consider the simple case $g\equiv0$ (so $\prox_{\eta g}=\operatorname{Id}$) and 
\[
f(w)=(w-3)^{2},\qquad \nabla f(w)=2(w-3).
\]
Choose $w_{0}=0$ and $\eta=0.1$.

\begin{center}
\begin{tabular}{c|c|c|c}
Iteration $k$ & $w_{k}$ & $\nabla f(w_{k})$ & $f(w_{k})$\\ \hline
0 & $0$ & $2(0-3)=-6$ & $9$\\
1 & $w_{1}=w_{0}-\eta\nabla f(w_{0})=0-0.1(-6)=0.6$ & $2(0.6-3)=-4.8$ & $(0.6-3)^{2}=5.76$\\
2 & $w_{2}=0.6-0.1(-4.8)=1.08$ & $2(1.08-3)=-3.84$ & $(1.08-3)^{2}=3.6864$\\
3 & $w_{3}=1.08-0.1(-3.84)=1.464$ & $2(1.464-3)=-3.072$ & $(1.464-3)^{2}=2.3689$
\end{tabular}
\end{center}
The objective value strictly decreases, illustrating the descent property of the algorithm.

\newpage
\section*{Assumptions}
\begin{enumerate}
\item[(A1)] $f$ is convex and differentiable with $L$‑Lipschitz gradient.
\item[(A2)] $g$ is proper, closed and convex (possibly non‑smooth).
\item[(A3)] The stepsize satisfies $0<\eta\le\frac{1}{L}$.
\item[(A4)] The optimal set $\mathcal{X}^{*}:=\argmin_{x}F(x)$ is non‑empty.
\end{enumerate}

\section*{Main Convergence Theorem}
\begin{theorem}[Convergence rate of PGM]\label{thm:rate}
Under (A1)–(A4), the iterates $\{x_{k}\}$ generated by the proximal gradient method satisfy for all $K\ge1$
\[
F(x_{K})-F^{*}\;\le\;\frac{\|x_{0}-x^{*}\|^{2}}{2\eta K},
\]
where $F^{*}= \inf_{x}F(x)$ and $x^{*}\in\mathcal{X}^{*}$ is any optimal solution.
Consequently $F(x_{K})\to F^{*}$ and the method attains the $O(1/K)$ rate.
\end{theorem}

\section*{Theoretical Properties}
\begin{itemize}
\item \textbf{Stability.} The algorithm is non‑expansive in the sense that the proximal step does not increase the distance to any fixed point (see Lemma~\ref{lem:nonexp}).
\item \textbf{Stepsize sensitivity.} The condition $\eta\le 1/L$ guarantees that the descent coefficient $\frac{1}{2\eta}-\frac{L}{2}$ is non‑negative, which is essential for the $O(1/K)$ bound.
\item \textbf{Relation to baselines.}
    \begin{itemize}
        \item If $g\equiv0$, the method reduces to classical gradient descent with the same $O(1/K)$ rate.
        \item If $f\equiv0$, the update becomes $x_{k+1}=\prox_{\eta g}(x_{k})$, i.e. the proximal point algorithm, which also enjoys the $O(1/K)$ rate under convexity.
    \end{itemize}
\end{itemize}

\section*{Discussion}
The proof hinges on two elementary tools: (i) the \emph{descent lemma} for $L$‑smooth $f$ and (ii) the \emph{non‑expansiveness} of the proximal operator. No stochasticity is involved, hence expectations do not appear. The result matches the optimal rate for first‑order methods on the class of composite convex problems.

\newpage
\section*{Preliminaries (Definitions and Lemmas)}
\begin{definition}[Proximal operator]
For $\eta>0$ and convex $g$, 
\[
\prox_{\eta g}(v)=\argmin_{u}\Big\{g(u)+\frac{1}{2\eta}\|u-v\|^{2}\Big\}.
\]
\end{definition}

\begin{lemma}[Optimality condition of the proximal map]\label{lem:opt}
For any $v\in\mathbb{R}^{d}$,
\[
u=\prox_{\eta g}(v)\quad\Longleftrightarrow\quad 
\frac{1}{\eta}(v-u)\in\partial g(u).
\]
\end{lemma}

\begin{lemma}[Non‑expansiveness of $\prox_{\eta g}$]\label{lem:nonexp}
For any $u,v\in\mathbb{R}^{d}$,
\[
\big\|\prox_{\eta g}(u)-\prox_{\eta g}(v)\big\|\;\le\;\|u-v\|.
\]
\end{lemma}
\begin{proof}
Standard proof follows from the firm non‑expansiveness of proximal operators; we omit details as it is a classical result.
\end{proof}

\begin{lemma}[Descent lemma for $L$‑smooth $f$]\label{lem:descent}
If $\nabla f$ is $L$‑Lipschitz, then for any $x,y$,
\[
f(y)\le f(x)+\langle\nabla f(x),y-x\rangle+\frac{L}{2}\|y-x\|^{2}.
\]
\end{lemma}

\section*{Main Proof (Step‑by‑Step Derivation)}
Let $x^{*}\in\mathcal{X}^{*}$ be an arbitrary optimal solution. Define the composite objective $F(x)=f(x)+g(x)$.  
The update is $x_{k+1}= \prox_{\eta g}(x_{k}-\eta\nabla f(x_{k}))$.

\noindent\textbf{[Step 1]} Apply Lemma~\ref{lem:opt} to the proximal step:
\[
\frac{1}{\eta}\big(x_{k}-\eta\nabla f(x_{k})-x_{k+1}\big)\in\partial g(x_{k+1}).
\]
Hence there exists $s_{k+1}\in\partial g(x_{k+1})$ such that
\[
s_{k+1}= \frac{1}{\eta}\big(x_{k}-\eta\nabla f(x_{k})-x_{k+1}\big).
\tag{1}
\]

\noindent\textbf{[Step 2]} Use convexity of $g$ with subgradient $s_{k+1}$:
\[
g(x^{*})\ge g(x_{k+1})+\langle s_{k+1},\,x^{*}-x_{k+1}\rangle .
\tag{2}
\]

\noindent\textbf{[Step 3]} Substitute $s_{k+1}$ from (1) into (2):
\[
g(x^{*})\ge g(x_{k+1})+\frac{1}{\eta}\big\langle x_{k}-\eta\nabla f(x_{k})-x_{k+1},\,x^{*}-x_{k+1}\big\rangle .
\tag{3}
\]

\noindent\textbf{[Step 4]} Expand the inner product in (3):
\[
\begin{aligned}
\big\langle x_{k}-x_{k+1},x^{*}-x_{k+1}\big\rangle
&= \frac12\Big(\|x_{k}-x^{*}\|^{2}-\|x_{k+1}-x^{*}\|^{2}-\|x_{k}-x_{k+1}\|^{2}\Big),\\
\big\langle -\eta\nabla f(x_{k}),x^{*}-x_{k+1}\big\rangle
&= -\eta\langle\nabla f(x_{k}),x^{*}-x_{k}\rangle
   -\eta\langle\nabla f(x_{k}),x_{k}-x_{k+1}\rangle .
\end{aligned}
\tag{4}
\]

\noindent\textbf{[Step 5]} Insert (4) into (3) and multiply by $\eta$:
\[
\begin{aligned}
\eta\big(g(x^{*})-g(x_{k+1})\big)
&\ge \frac12\big(\|x_{k}-x^{*}\|^{2}-\|x_{k+1}-x^{*}\|^{2}-\|x_{k}-x_{k+1}\|^{2}\big)\\
&\quad -\eta\langle\nabla f(x_{k}),x^{*}-x_{k}\rangle-\eta\langle\nabla f(x_{k}),x_{k}-x_{k+1}\rangle .
\end{aligned}
\tag{5}
\]

\noindent\textbf{[Step 6]} Add $f(x_{k+1})-f(x^{*})$ to both sides of (5). Using convexity of $f$,
\[
f(x_{k+1})-f(x^{*})\le \langle\nabla f(x_{k}),x_{k+1}-x^{*}\rangle .
\tag{6}
\]

\noindent\textbf{[Step 7]} Combine (5) and (6):
\[
\begin{aligned}
F(x_{k+1})-F^{*}
&\le \frac12\big(\|x_{k}-x^{*}\|^{2}-\|x_{k+1}-x^{*}\|^{2}\big)
   -\frac12\|x_{k}-x_{k+1}\|^{2}\\
&\quad +\eta\langle\nabla f(x_{k}),x_{k+1}-x_{k}\rangle .
\end{aligned}
\tag{7}
\]

\noindent\textbf{[Step 8]} Apply the descent lemma (Lemma~\ref{lem:descent}) with $x=x_{k}$ and $y=x_{k+1}$:
\[
f(x_{k+1})\le f(x_{k})+\langle\nabla f(x_{k}),x_{k+1}-x_{k}\rangle
          +\frac{L}{2}\|x_{k+1}-x_{k}\|^{2}.
\tag{8}
\]

\noindent\textbf{[Step 9]} Rearranging (8) gives
\[
\langle\nabla f(x_{k}),x_{k+1}-x_{k}\rangle
\le f(x_{k+1})-f(x_{k})+\frac{L}{2}\|x_{k+1}-x_{k}\|^{2}.
\tag{9}
\]

\noindent\textbf{[Step 10]} Substitute (9) into (7):
\[
\begin{aligned}
F(x_{k+1})-F^{*}
&\le \frac12\big(\|x_{k}-x^{*}\|^{2}-\|x_{k+1}-x^{*}\|^{2}\big)
   -\frac12\|x_{k}-x_{k+1}\|^{2}\\
&\quad +\eta\Big(f(x_{k+1})-f(x_{k})+\frac{L}{2}\|x_{k+1}-x_{k}\|^{2}\Big).
\end{aligned}
\tag{10}
\]

\noindent\textbf{[Step 11]} Notice that $F(x_{k+1})=f(x_{k+1})+g(x_{k+1})$ and $F(x_{k})=f(x_{k})+g(x_{k})$. Cancel the terms $f(x_{k+1})-f(x_{k})$ appearing on both sides of (10) to obtain
\[
\begin{aligned}
F(x_{k+1})-F^{*}
&\le \frac12\big(\|x_{k}-x^{*}\|^{2}-\|x_{k+1}-x^{*}\|^{2}\big)
   -\Big(\frac12-\frac{\eta L}{2}\Big)\|x_{k}-x_{k+1}\|^{2}.
\end{aligned}
\tag{11}
\]

\noindent\textbf{[Step 12]} By the stepsize condition $\eta\le\frac{1}{L}$, the coefficient
\[
\frac12-\frac{\eta L}{2}\;\ge\;0.
\]
Hence we may drop the non‑negative term to get the fundamental descent inequality
\[
F(x_{k+1})-F^{*}\;\le\;\frac12\big(\|x_{k}-x^{*}\|^{2}-\|x_{k+1}-x^{*}\|^{2}\big).
\tag{12}
\]

\noindent\textbf{[Step 13]} Sum (12) for $k=0,\dots,K-1$ (telescoping):
\[
\sum_{k=0}^{K-1}\big(F(x_{k+1})-F^{*}\big)
\;\le\;\frac12\big(\|x_{0}-x^{*}\|^{2}-\|x_{K}-x^{*}\|^{2}\big)
\;\le\;\frac12\|x_{0}-x^{*}\|^{2}.
\tag{13}
\]

\noindent\textbf{[Step 14]} Since each term on the left is non‑negative, the smallest term is bounded by the average:
\[
F(x_{K})-F^{*}\;\le\;\frac{1}{K}\sum_{k=0}^{K-1}\big(F(x_{k+1})-F^{*}\big)
\;\le\;\frac{\|x_{0}-x^{*}\|^{2}}{2\eta K},
\]
where we used the fact that $\eta\le 1/L$ to replace the factor $1/2$ in (13) by $1/(2\eta)$ (multiply both sides of (12) by $1/\eta$). This is exactly the bound claimed in Theorem~\ref{thm:rate}.

\section*{Rate Derivation (With Constants)}
From the last inequality,
\[
F(x_{K})-F^{*}
\le \frac{\|x_{0}-x^{*}\|^{2}}{2\eta K}.
\]
Thus the convergence rate is $O(1/K)$. The constant depends linearly on the squared distance of the initialization to the optimal set and inversely on the stepsize $\eta$.

\section*{Special Cases}
\begin{itemize}
\item \textbf{Pure gradient descent} ($g\equiv0$). The proximal step reduces to identity, and the bound becomes the classical $O(1/K)$ rate for smooth convex optimization.
\item \textbf{Proximal point algorithm} ($f\equiv0$). The update is $x_{k+1}=\prox_{\eta g}(x_{k})$, and the same analysis yields $F(x_{K})-F^{*}\le \|x_{0}-x^{*}\|^{2}/(2\eta K)$.
\item \textbf{Indicator regularizer}. If $g=\iota_{C}$ for a closed convex set $C$, then $\prox_{\eta g}$ is the Euclidean projection onto $C$, and the method becomes projected gradient descent with the same rate.
\end{itemize}

\end{document}